%%%%%%%%%%%%%%%%%
% This is an sample CV template created using altacv.cls
% (v1.1.4, 27 July 2018) written by LianTze Lim (liantze@gmail.com). Now compiles with pdfLaTeX, XeLaTeX and LuaLaTeX.
% 
%% It may be distributed and/or modified under the
%% conditions of the LaTeX Project Public License, either version 1.3
%% of this license or (at your option) any later version.
%% The latest version of this license is in
%%    http://www.latex-project.org/lppl.txt
%% and version 1.3 or later is part of all distributions of LaTeX
%% version 2003/12/01 or later.
%%%%%%%%%%%%%%%%

%% If you need to pass whatever options to xcolor
\PassOptionsToPackage{dvipsnames}{xcolor}

%% Disable pdf/A
%\PassOptionsToPackage{disable}{pdfx}


%% If you are using \orcid or academicons
%% icons, make sure you have the academicons 
%% option here, and compile with XeLaTeX
%% or LuaLaTeX.
\documentclass[10pt,a4paper,withhyper]{altacv}

%% Use the "normalphoto" option if you want a normal photo instead of cropped to a circle
% \documentclass[10pt,a4paper,normalphoto]{altacv}

%\documentclass[10pt,a4paper]{altacv}
%% AltaCV uses the fontawesome and academicon fonts
%% and packages. 
%% See texdoc.net/pkg/fontawecome and http://texdoc.net/pkg/academicons for full list of symbols.
%% 
%% Compile with LuaLaTeX for best results. If you
%% want to use XeLaTeX, you may need to install
%% Academicons.ttf in your operating system's font 
%% folder.


% Change the page layout if you need to
\geometry{
  left=0.5cm,
  right=7.5cm,
  marginparwidth=6.8cm,
  marginparsep=0.3cm,
  top=0.5cm,
  bottom=0.5cm,
  footskip=2\baselineskip
}

% Change the font if you want to.

% If using pdflatex:
%\usepackage[T1]{fontenc}
%\usepackage[utf8]{inputenc}
%\usepackage[default]{lato}

% If using xelatex or lualatex:
\setmainfont{Lato}

% Change the colours if you want to
\definecolor{Navy}{HTML}{000080}
\definecolor{SlateGrey}{HTML}{2E2E2E}
\definecolor{LightGrey}{HTML}{666666}
\colorlet{heading}{Navy}
\colorlet{accent}{Navy}
\colorlet{emphasis}{SlateGrey}
\colorlet{body}{LightGrey}


% Change the bullets for itemize and rating marker
% for \cvskill if you want to
\renewcommand{\itemmarker}{{\small\textbullet}}
\renewcommand{\ratingmarker}{\faCircle}

\renewcommand{\divider}{%
  \vspace{-0.15\baselineskip}%
  \textcolor{body!25}{\hdashrule{\linewidth}{0.4pt}{0.8ex}}%
  \vspace{2pt}%
}



%% sample.bib contains your publications
%\usepackage[
%  backend=biber,
%  bibencoding=utf8
%]{biblatex}
%\addbibresource{\jobname.bib}
%\addbibresource{mybib2.bib}

% \usepackage[colorlinks]{hyperref}

%%%%%%%%%%%%%%%%%
% Resume content optimized for ATS + Dynamic Radiant Technique (DRT)
%%%%%%%%%%%%%%%%

\begin{document}

\name{Kyle Dawson, PhD}
\tagline{PhD Scientist \& Cloud Solution Builder | Satellite Remote Sensing, Aerospace \& Machine Learning}
\personalinfo{%
  \email{hello@kyledawson.me}
  \phone{+1 (618)535-9822}
  \homepage{kyledawson.me}
  \linkedin{kylewdawson}
  \github{kyledawson}
  \orcid{0000-0003-3175-0456}
}

\begin{fullwidth}
\makecvheader
\end{fullwidth}

%======================
% EXPERIENCE
%======================

\cvsectionwithicon[page1sidebar]{\faIcon{briefcase}}{Experience}

\cvevent{Lead Data Scientist — Energy}{Centerpoint Energy}{Apr 2026 -- Present}{Houston, TX}
\begin{itemize}
  \item lead a new geospatial intelligence (GEOINT) team, setting technical direction and architecture standards for satellite-, LiDAR-, and imagery-driven systems
  \item architected and deployed cloud-native, containerized services (Cloud Run, Docker, CI/CD, BigQuery) that ingest and process large-scale geospatial and remote-sensing data through scalable, reliable pipelines
  \item built a GPU-accelerated 3D computer-vision pipeline that maps infrastructure by aligning engineering CAD models to LiDAR point clouds
  \item prototyped and deployed a multi-model generative-AI application integrating managed foundation-model services (Vertex AI, Cohere) with a retrieval benchmark for document Q\&A and image analysis
\end{itemize}

\divider

\cvevent{Data Scientist → Senior Data Scientist — Climate Tech}{SeekOps Inc.}{Mar 2022 -- Apr 2026}{Austin, TX}
\begin{itemize}
\item designed, deployed, and delivered customer-facing cloud analytics platforms (Plotly Dash on Azure) that turned complex satellite- and drone-based sensing capabilities into productized solutions, directly securing contracts of \$50--150k
\item served as customer-facing technical lead across 10+ countries, presenting solutions, architecture, and technical strategy to executives, enterprise customers, and international conference audiences
\item executed proofs of concept and authored technical documentation and uncertainty assessments supporting multi-year enterprise contracts
\end{itemize}

\divider

\cvevent{Data Scientist — Ag Tech}{Bayer CropScience}{Jan 2020 -- Mar 2022}{St. Louis, MO}
\begin{itemize}
\item built end-to-end cloud data workflows on AWS (S3) combining relational databases, geospatial data, and aerial imagery into analytics-ready datasets
\item applied machine learning and computer vision (Detectron2) to UAV- and aircraft-derived imagery for large-scale digital phenotyping
\item delivered interactive decision-support applications to internal stakeholders and business teams
\end{itemize}

\divider

\cvevent{Postdoctoral Fellow — Aerospace \& Remote Sensing}{NASA Langley Research Center}{Feb 2018 -- Jan 2020}{Hampton, VA}
\begin{itemize}
\item developed retrieval algorithms for airborne and satellite remote-sensing instruments (combined High Spectral Resolution and Raman LiDAR) supporting NASA Earth-observation and atmospheric-science missions
\item delivered peer-reviewed results and instrument data tooling for the aerospace remote-sensing research community
\item prior NASA appointments at Langley (LARSS, HSRL) and Ames (active/passive remote-sensing algorithm comparison), spanning three NASA research centers
\end{itemize}

%======================
% SELECTED PROJECT
%======================

\cvsectionwithicon{\faIcon{satellite}}{Selected Project}
\cvevent{NASA HSRL Data Toolkit (hsrltool)}{Open-source scientific software}{}{}{A graphical toolkit for reading, inspecting, and visualizing NASA High Spectral Resolution Lidar (HSRL) instrument data (HDF5), used by the atmospheric remote-sensing community.}~{\href{https://github.com/kyledawson/hsrltool}{\faIcon{external-link-alt}}}

%======================
% EDUCATION
%======================

\cvsectionwithicon{\faIcon{graduation-cap}}{Education}
{\small
\cvevent{PhD, Atmospheric Science}{North Carolina State University}{Graduated 2017}{Raleigh, NC}{Contributions to the Understanding of Aerosol Microphysics Towards Improving the Assessment of Climate Radiative Forcing}~{\href{https://repository.lib.ncsu.edu/handle/1840.20/34316}{\faIcon{external-link-alt}}}
}
\end{document}
